\begin{lemma}[A disjoint swap preserves later elementary-prefix missing sets]
Let \(\mu\) be a proper partial coloring of a finite loopless
multigraph, let \(\alpha\ne\beta\), and let exact tree-sequences
\(B\preccurlyeq X\preccurlyeq Z\) be related by
\noderef{TreeSequenceInitialSegment}. Assume \(V(Z)\) is
\noderef{ColoredVertexSetElementary} under \(\mu\), both colors
belong to \(\noderef{TreeMissingColors}(\mu,B)\), and the
\noderef{TwoColorComponent} \(Q\) rooted at an arbitrary incidence
node \(z\) avoids every vertex-node of \(X\).
For \(\sigma=\noderef{TwoColorComponentSwap}(G,\mu,\alpha,\beta,z)\),
every edge incident with \(X\) keeps its color, and
\(\noderef{MissingColors}(G,\sigma,v)=
\noderef{MissingColors}(G,\mu,v)\) for every \(v\in V(Z)\).
\end{lemma}
\begin{proof}
An affected-color edge-node is adjacent to each of its endpoints in
\noderef{TwoColorIncidenceGraph}. Adjacent nodes belong to the same
\noderef{TwoColorComponent}. Thus an affected-color edge of \(Q\)
cannot be incident with \(X\), and the swap fixes every incident edge:
off-color values are fixed even on edge-nodes in \(Q\).
For missing sets on \(Z\), vertices outside \(Q\) have unchanged
missing sets by \noderef{TwoColorComponentSwapMissingBehavior}.
Fix a vertex \(v\in V(Z)\) lying in \(Q\). For either affected
color \(c\), choose its missing witness \(w\in V(B)\).
\noderef{TreeSequenceInitialSegmentVertexSetSubset} puts \(w\)
in \(X\) and \(Z\). Hence \(w\notin Q\) and \(w\ne v\);
elementarity implies that \(c\) is not missing at \(v\).
This reasoning applies separately to both colors and does not require
their missing witnesses to be distinct.
There is consequently an old incident edge of each color at \(v\).
Both edge-nodes lie in \(Q\), so they exchange colors under the swap.
Both affected colors are therefore also present at \(v\) afterwards.
For each other color, its missing predicate is unchanged by
\noderef{TwoColorComponentSwapMissingBehavior}. Checking membership
color by color proves equality of the missing sets at \(v\).
\end{proof}
